
En
effet, on a d\'ej\`a not\'e l'inclusion \(\Phi(G'^\tau)\supset G^\tau\). Pour l'inclusion en sens inverse, on note que \(g\in G(K)\) est dans \(G^\tau(K)\) si et seulement si la suite des \(ng\) (\(n\in\mathbb N\)) est quasi-compacte\,; comme \(\Phi_K\) est quasi-compact, il suffit pour ceci que la suite des \(\Phi(ng)=n\Phi(g)\) ait la m\^eme propri\'et\'e\,; d'o\`u l'assertion.

\textbf{Lemme 4.3.} \emph{Soient \(X\), \(Y\) propres sur un sch\'ema \(S\) et \(f:X\to Y\) un morphisme surjectif (resp. le morphisme d'inclusion du sch\'ema des z\'eros d'une section \(\varphi\) d'un faisceau ample \({O}_Y(1)\)). Alors}
\[
  (f^*)^{-1}\underline{\Pic}^\tau_{X/S}=\underline{\Pic}^\tau_{Y/S}.
\]

En effet, l'assertion est cons\'equence formelle de 4.2, (XII 1.2), 3.1, et 3.5 (resp. 3.8).

\textbf{4.4.} Soient \(k\) un corps alg\'ebriquement clos, \(X\) un \(k\)-sch\'ema propre, et \(L\) un faisceau inversible sur \(X\). On dit que \(L\) est \uline{alg\'ebriquement \'equivalent} \`a \(0\) (resp. \(\tau\)-\uline{\'equivalent} \`a \(0\)) s'il existe un \(k\)-sch\'ema connexe de type fini \(T\), un faisceau inversible \(M\) sur \(X_T=X\times_k T\), et deux \(k\)-points \(t,t_0\) de \(T\) tels que \(L\simeq M_t\) et \({O}_X\simeq M_{t_0}\) (resp. s'il existe un entier non nul \(n\) tel que \(L^{\otimes n}\) soit alg\'ebriquement \'equivalent \`a \(0\)), ou ce qui revient facilement au m\^eme, si le \(k\)-point \(\lambda\) de \(\underline{\Pic}_{X/k}\) qui correspond \`a \(L\), est dans \(\underline{\Pic}^0_{X/k}(k)\) (resp. dans \(\underline{\Pic}^\tau_{X/k}(k)\)).

Rappelons aussi que \(L\) est dit \uline{num\'eriquement \'equivalent} \`a \(0\) si pour toute courbe int\`egre ferm\'ee \(Y\) dans \(X\), le nombre d'intersection \(\langle c_1(L)\cdot Y\rangle\) est \(0\) (voir 2.1 ou X 4.5).

